% Luc Destombe --- licence CC BY-NC-SA 4.0
% https://creativecommons.org/licenses/by-nc-sa/4.0/deed.fr
% Fichier autonome : compiler avec pdflatex, deux fois (signets, nombre de pages).
\documentclass[11pt,a4paper]{article}
\makeatletter
\newif\ifeleve
\elevefalse
\elevetrue

\newif\iflycee@palatino
\lycee@palatinotrue

\RequirePackage[utf8]{inputenc}
\RequirePackage[T1]{fontenc}
\RequirePackage[french]{babel}
\RequirePackage{etoolbox}

\newcommand{\ShorthandsOff}{\@ifundefined{shorthandoff}{}{\shorthandoff{;:!?}}}
\newcommand{\ShorthandsOn}{\@ifundefined{shorthandon}{}{\shorthandon{;:!?}}}
\ShorthandsOff
\AtBeginDocument{\ShorthandsOn}
\AtBeginEnvironment{tikzpicture}{\ShorthandsOff}

\RequirePackage{geometry}
\RequirePackage{fancyhdr}
\RequirePackage{lastpage}
\RequirePackage{multicol}
\RequirePackage{array}
\RequirePackage{enumitem}
\RequirePackage{xcolor}
\RequirePackage{needspace}

\RequirePackage{amsmath,amssymb}
\RequirePackage{mathpazo}
\DeclareMathSizes{10.95}{10.95}{8}{5.5}
\begingroup\catcode`\_=\active \catcode`\^=\active
\gdef_#1{\sb{\scriptscriptstyle #1}}
\gdef^#1{\sp{\scriptscriptstyle\lycee@moinsepais #1}}
\endgroup
\newcommand\lycee@moins{\mathbin{%
  \setbox\z@\hbox{$\m@th\scriptscriptstyle\lycee@moinsorig$}%
  \dimen@=.55\wd\z@ \advance\dimen@-.5pt
  \hbox to.75\wd\z@{\hss$\m@th\scriptscriptstyle\vcenter{\hbox{%
    \tikz\draw[line width=.5pt,line cap=round](0,0)--(\the\dimen@,0);}}$\hss}}}
\begingroup\lccode`\~=`\- \lowercase{\endgroup\let~}\lycee@moins
\newcommand\lycee@moinsepais{\mathcode`\-="8000 }
\AtBeginDocument{\mathchardef\lycee@moinsorig=\mathcode`\-\relax}
\AtBeginDocument{\catcode`\_=12 \mathcode`\_="8000
  \catcode`\^=12 \mathcode`\^="8000 }
\newcommand\lycee@exposants{%
  \fontdimen13\textfont2=.48\fontdimen6\textfont2
  \fontdimen14\textfont2=.48\fontdimen6\textfont2
  \fontdimen15\textfont2=.48\fontdimen6\textfont2
  \fontdimen13\scriptfont2=.48\fontdimen6\scriptfont2
  \fontdimen14\scriptfont2=.48\fontdimen6\scriptfont2
  \fontdimen15\scriptfont2=.48\fontdimen6\scriptfont2}
\every@math@size\expandafter{\the\every@math@size\lycee@exposants}
\newlength{\retraitcalculs}
\setlength{\retraitcalculs}{2em}

\RequirePackage{tikz}
\usetikzlibrary{arrows.meta,calc,babel,shapes.misc}
\RequirePackage[most]{tcolorbox}
\tcbuselibrary{skins,breakable}

\definecolor{coulDef}{HTML}{1F4E79}
\definecolor{coulProp}{HTML}{7030A0}
\definecolor{coulRem}{HTML}{C55A11}
\definecolor{coulEx}{HTML}{2E7D32}
\definecolor{coulExo}{HTML}{B03A2E}
\definecolor{coulPart}{HTML}{1F4E79}
\definecolor{coulMeth}{HTML}{1B6E4F}
\definecolor{coulCourbe}{HTML}{0B5ED7}
\definecolor{coulCourbeB}{HTML}{C00000}
\definecolor{coulCourbeC}{HTML}{2E7D32}
\definecolor{coulGrille}{HTML}{8C8C8C}
\definecolor{coulRep}{HTML}{1B6E4F}
\definecolor{coulBar}{HTML}{2E7D32}

\newcommand{\R}{\mathbb{R}}
\newcommand{\Rs}{\mathbb{R}^{*}}
\renewcommand{\le}{\leqslant}
\renewcommand{\ge}{\geqslant}
\renewcommand{\approx}{\simeq}
\newcommand{\vect}[1]{\overrightarrow{#1}}
\newcommand{\norme}[1]{\left\lVert #1 \right\rVert}
\newcommand{\intervalle}[2]{\left[#1\,;#2\right]}
\RequirePackage{cancel}
\renewcommand{\CancelColor}{\color{coulExo}}
\newcommand{\fois}[1]{\times{\color{coulExo}#1}}
\newcommand{\barre}[1]{\cancel{#1}}

\tikzset{grille/.style={coulGrille,line width=0.4pt}}
\newcommand{\quadrillage}[4]{\draw[grille,step=1] (#1,#2) grid (#3,#4);}
\tikzset{
  croix/.style={cross out,draw,line width=0.6pt,inner sep=0pt,minimum size=4pt},
  croixdroite/.style={croix,rotate=45,line width=0.8pt},
}
\newcommand{\pt}[3]{\path (#1) node[croix]{} node[#2,font=\small,inner sep=2.5pt]{$#3$};}

\newcommand{\lycee@repere}[4]{%
  \draw[grille,step=1] (#1,#2) grid (#3,#4);
  \draw[-{Stealth[length=2mm]},thick] (#1,0)--({#3+0.4},0) node[below left=-1pt]{$x$};
  \draw[-{Stealth[length=2mm]},thick] (0,#2)--(0,{#4+0.4}) node[below left=-1pt]{$y$};
  \node[below left,font=\scriptsize,inner sep=1.5pt] at (0,0) {$0$};}
\newcommand{\lycee@gradx}[1]{\foreach \k/\l in {#1}{%
  \draw (\k,0.07)--(\k,-0.07) node[below,font=\scriptsize,inner sep=1.5pt]{$\l$};}}
\newcommand{\lycee@grady}[1]{\foreach \k/\l in {#1}{%
  \draw (0.07,\k)--(-0.07,\k) node[left,font=\scriptsize,inner sep=1.5pt]{$\l$};}}
\newcommand{\lycee@intff}[2]{\left[#1\,;#2\right]}
\newcommand{\lycee@intfo}[2]{\left[#1\,;#2\right[}
\newcommand{\lycee@intof}[2]{\left]#1\,;#2\right]}
\newcommand{\lycee@intoo}[2]{\left]#1\,;#2\right[}
\newcommand{\lycee@ens}[1]{\left\{#1\right\}}
\AtBeginDocument{%
  \providecommand{\repere}{\lycee@repere}%
  \providecommand{\gradx}{\lycee@gradx}%
  \providecommand{\grady}{\lycee@grady}%
  \providecommand{\Cf}{\mathcal{C}\sb{\scriptscriptstyle f}}%
  \providecommand{\Cg}{\mathcal{C}\sb{\scriptscriptstyle g}}%
  \providecommand{\intff}{\lycee@intff}%
  \providecommand{\intfo}{\lycee@intfo}%
  \providecommand{\intof}{\lycee@intof}%
  \providecommand{\intoo}{\lycee@intoo}%
  \providecommand{\ens}{\lycee@ens}}

\newlist{questions}{enumerate}{3}
\setlist[questions,1]{label=\arabic*\textdegree),leftmargin=*,itemsep=2pt,topsep=3pt}
\setlist[questions,2]{label=\alph*),leftmargin=*,itemsep=1pt,topsep=2pt}
\setlist[questions,3]{label=\roman*),leftmargin=*,itemsep=1pt,topsep=1pt}
\setlist[itemize]{label=\textbullet,leftmargin=*,itemsep=2pt,topsep=3pt}

\newcommand{\besoinplace}[1]{\par\penalty\@M\begingroup
  \dimen@=#1\relax\dimen@ii=\pagegoal\advance\dimen@ii-\pagetotal
  \ifdim\dimen@>\dimen@ii\ifdim\dimen@ii>\z@\vfil\fi\break\fi\endgroup}

\newcommand{\lycee@pied}{\thepage/\pageref{LastPage}}
\newcommand{\entete}[2]{%
  \pagestyle{fancy}\fancyhf{}%
  \fancyhead[L]{\footnotesize\color{coulPart}\textbf{#1}}%
  \fancyhead[R]{\footnotesize\color{coulPart}\textbf{#2}}%
  \fancyfoot[C]{\footnotesize\lycee@pied}%
  \renewcommand{\headrulewidth}{0.4pt}%
  \renewcommand{\headrule}{\hbox to\headwidth{%
    \color{coulPart!40}\leaders\hrule height \headrulewidth\hfill}}}

\RequirePackage{ccicons}
\newcommand{\lycee@licence}{{\scriptsize\color{gray}%
  \copyright\ Luc Destombe --- \ccbyncsa\ CC BY-NC-SA 4.0}}
\newif\iflycee@licence \lycee@licencetrue
\newcommand{\sanslicence}{\lycee@licencefalse}
\AtBeginDocument{\iflycee@licence\AddToHookNext{shipout/foreground}{%
  \put(0,0){\rlap{\kern\dimexpr1in+\hoffset+\oddsidemargin+\textwidth\relax
    \llap{\raisebox{-\dimexpr1in+\voffset+\topmargin+\headheight+\headsep
      +\textheight+\footskip\relax}{\lycee@licence}}}}}\fi}

\newcommand{\formulesserrees}{%
  \setlength{\abovedisplayskip}{3pt}\setlength{\belowdisplayskip}{3pt}%
  \setlength{\abovedisplayshortskip}{2pt}\setlength{\belowdisplayshortskip}{3pt}}

\setlength{\parindent}{0pt}

\RequirePackage[implicit=false,hidelinks,pdfencoding=auto,psdextra,bookmarksopen,
  bookmarksopenlevel=1,pdfcreator={},pdfproducer={}]{hyperref}
\RequirePackage{bookmark}
\hypersetup{pdfauthor={Luc Destombe},pdfkeywords={CC BY-NC-SA 4.0}}
\newcounter{lycee@signet}
\newcommand{\signet}[3][]{\stepcounter{lycee@signet}%
  \edef\lycee@dest{\if\relax\detokenize{#1}\relax
    signet.\arabic{lycee@signet}\else#1\fi}%
  \hypertarget{\lycee@dest}{}\bookmark[level=#2,dest=\lycee@dest]{#3}}
\pdfstringdefDisableCommands{%
  \def\R{R}\def\vect#1{#1}\def\Cf{Cf}\def\Cg{Cg}\def\fois#1{×#1}%
  \def\barre#1{#1}\def\intervalle#1#2{[#1;#2]}\def\intff#1#2{[#1;#2]}%
  \def\textsuperscript#1{#1}\def\dfrac#1#2{#1/#2}\def\frac#1#2{#1/#2}%
  \def\vec#1{#1⃗}}%
\geometry{top=2cm,bottom=1cm,left=1cm,right=1cm,headheight=14pt,footskip=0.6cm}

\RequirePackage{pgfplots}
\pgfplotsset{compat=1.18}
\RequirePackage{tkz-tab}

\newcounter{partiecours}
\renewcommand{\thepartiecours}{\Roman{partiecours}}
\newcommand{\partiecours}[1]{%
  \refstepcounter{partiecours}%
  \Needspace*{8\baselineskip}%
  \bigskip
  \noindent\begin{tcolorbox}[enhanced,colback=coulPart,colframe=coulPart,
      boxrule=0pt,arc=2pt,left=8pt,right=8pt,top=4pt,bottom=4pt,
      before skip=6pt,after skip=10pt]
    \color{white}\bfseries\large \thepartiecours{} --- #1%
    \signet[partie-\arabic{partiecours}]{1}{\thepartiecours{} --- #1}
  \end{tcolorbox}%
  \addcontentsline{toc}{section}{\thepartiecours{} --- #1}%
}

\newtcolorbox{boitecours}[3][]{%
  enhanced,breakable,
  colback=white,colframe=#2,coltitle=white,
  fonttitle=\bfseries,
  attach boxed title to top left={xshift=8pt,yshift=-\tcboxedtitleheight/2},
  boxed title style={colback=#2,arc=2pt,boxrule=0pt},
  boxrule=0.9pt,arc=2pt,
  left=8pt,right=8pt,top=8pt,bottom=6pt,
  title={#3},#1}

\newcounter{methode}
\newenvironment{definitionbox}[1][Définition]%
  {\begin{boitecours}[phantom={\signet{2}{#1}}]{coulDef}{#1}}{\end{boitecours}}
\newenvironment{proprietebox}[1][Propriété]%
  {\begin{boitecours}[phantom={\signet{2}{#1}}]{coulProp}{#1}}{\end{boitecours}}
\newenvironment{methodebox}[1][Méthode]{\stepcounter{methode}%
  \begin{boitecours}[phantom={\signet[methode-\arabic{methode}]{2}{#1}}]{coulMeth}{#1}}%
  {\end{boitecours}}

\newcommand{\etape}[1]{\par\smallskip\textbf{\color{coulMeth}#1}\par\nobreak\smallskip}
\newcommand{\enonce}[1]{\emph{#1}\par\smallskip}
\newcommand{\reponse}[1]{\par\smallskip
  {\footnotesize\itshape\color{black!55}Réponses : #1}}
\newcommand{\demo}[1]{\textbf{\color{coulMeth}Démonstration.} #1}
\newcommand{\afaire}[1]{%
  \par\smallskip
  \noindent{\small\color{coulExo}$\blacktriangleright$\ \textbf{À faire :}
    fiche d'exercices du chapitre, #1.}\par\smallskip}

\newenvironment{calculs}{\setlength{\jot}{5pt}\@fleqntrue\@mathmargin\retraitcalculs
  \setlength{\abovedisplayskip}{4pt}\setlength{\belowdisplayskip}{4pt}%
  \setlength{\abovedisplayshortskip}{4pt}\setlength{\belowdisplayshortskip}{4pt}%
  \csname alignat*\endcsname{2}}%
  {\csname endalignat*\endcsname}
\newcommand{\rem}[1]{\text{\small\itshape\color{black!60}#1}}
\newcommand{\explique}[2]{\par\smallskip\noindent
  \begin{minipage}[t]{0.57\linewidth}\raggedright #1\end{minipage}\hfill
  \begin{minipage}[t]{0.40\linewidth}\raggedright\leavevmode
    {\small\itshape\color{black!60}#2}\end{minipage}\par}
\newcommand{\cas}[1]{\par\smallskip\noindent\textbf{\color{coulExo}#1}\nobreak}
\newcommand{\calc}[1]{\par\smallskip\textbf{\color{coulExo}#1}\par\nobreak}

\newtcolorbox{remarquebox}[1][Remarque]{%
  enhanced,breakable,colback=white,colframe=coulRem,
  boxrule=0pt,leftrule=2.5pt,arc=0pt,outer arc=0pt,
  left=8pt,right=6pt,top=5pt,bottom=5pt,
  before upper={\textbf{\color{coulRem}#1}\par\smallskip}}

\newtcolorbox{exemplebox}[1][Exemple]{%
  enhanced,breakable,colback=white,colframe=coulEx,
  boxrule=0pt,leftrule=2.5pt,arc=0pt,outer arc=0pt,
  left=8pt,right=6pt,top=5pt,bottom=5pt,
  before upper={\textbf{\color{coulEx}#1}\par\smallskip}}

\pgfplotsset{
  repere/.style={
    axis lines=middle,
    axis line style={-{Stealth[length=2.4mm]},thick},
    every axis x label/.style={at={(ticklabel* cs:1.0)},anchor=west},
    every axis y label/.style={at={(ticklabel* cs:1.0)},anchor=south},
    label style={font=\footnotesize},
    tick label style={font=\scriptsize},
    grid=both,
    major grid style={grille},
    minor tick num=0,
    tick align=inside,
    clip mode=individual,
  },
}
\tikzset{
  courbe/.style={coulCourbe,line width=1.1pt,smooth},
  courbeB/.style={coulCourbeB,line width=1.1pt,smooth},
  courbeC/.style={coulCourbeC,line width=1.1pt,smooth},
}

\newlist{etapes}{enumerate}{1}
\setlist[etapes]{label=\textbf{\arabic*.},leftmargin=*,itemsep=1pt,topsep=2pt}

\newlist{expressions}{enumerate}{1}
\setlist[expressions]{label={\Alph*\ $=$},leftmargin=*,itemsep=3pt,topsep=3pt}

\newcounter{exo}
\renewcommand{\theexo}{\ifnum\value{exo}<10 0\fi\arabic{exo}}
\newtcolorbox{exercice}[1][]{%
  enhanced,breakable,
  colback=white,colframe=coulExo,
  boxrule=0.9pt,arc=2pt,
  left=8pt,right=8pt,top=8pt,bottom=6pt,
  coltitle=white,fonttitle=\bfseries,
  attach boxed title to top left={xshift=8pt,yshift=-\tcboxedtitleheight/2},
  boxed title style={colback=coulExo,arc=2pt,boxrule=0pt},
  phantom={\signet{2}{Exercice \arabic{exo}}},
  title={Exercice \theexo\ifx\relax#1\relax\else{} --- #1\fi}}
\newenvironment{exo}[1][\relax]{\refstepcounter{exo}\begin{exercice}[#1]}{\end{exercice}}

  \RequirePackage{comment}
  \excludecomment{correction}
  \renewcommand{\reponse}[1]{}
  \newcommand{\autableau}{\hfill{\footnotesize\itshape\color{black!45}(au tableau)}}
  \newcommand{\etapeexemples}{%
    \par\smallskip\textbf{\color{coulMeth}Exemples}\autableau\par\nobreak\smallskip}
  \newcommand{\etapetableau}[1]{%
    \par\smallskip\textbf{\color{coulMeth}#1}\autableau\par\nobreak\smallskip}
  \newtcolorbox{exempletableau}[1][Exemple]{%
    enhanced,breakable,colback=white,colframe=coulEx,
    boxrule=0pt,leftrule=2.5pt,arc=0pt,outer arc=0pt,
    left=8pt,right=6pt,top=4pt,bottom=4pt,
    before skip=5pt,after skip=5pt,
    before upper={\textbf{\color{coulEx}#1}\autableau\par\smallskip}}

\RequirePackage{comment}
\excludecomment{correctionavj}

\newcommand{\coursEntete}{}
\newcommand{\coursNiveau}{}
\newcommand{\coursTitre}{}
\newcommand{\coursSurTitre}{}
\newcommand{\coursSousTitre}{}

\newcommand{\enteteCours}[1]{\renewcommand{\coursEntete}{#1}}
\newcommand{\chapitrecours}[2]{\enteteCours{Chapitre #1 --- #2}}
\newcommand{\niveaucours}[1]{\renewcommand{\coursNiveau}{#1}}
\newcommand{\titrecours}[3][]{%
  \renewcommand{\coursTitre}{#2}%
  \renewcommand{\coursSurTitre}{#1}%
  \renewcommand{\coursSousTitre}{#3}}

\pagestyle{fancy}
\fancyhf{}
\fancyhead[L]{\footnotesize\color{coulPart}\textbf{\coursEntete}}
\fancyhead[R]{\footnotesize\color{coulPart}\textbf{\coursNiveau}}
\fancyfoot[C]{\footnotesize\thepage}
\renewcommand{\headrulewidth}{0.4pt}
\renewcommand{\headrule}{\hbox to\headwidth{\color{coulPart!40}\leaders\hrule height \headrulewidth\hfill}}

\setlength{\parskip}{2pt}

\AtBeginDocument{%
  \begin{center}
    \begin{tcolorbox}[enhanced,width=0.72\linewidth,colback=white,colframe=coulPart,
        boxrule=1.6pt,arc=3pt,halign=center,top=10pt,bottom=10pt]
      \ifx\coursSurTitre\empty
        {\Huge\bfseries\color{coulPart} \coursTitre}\\[4pt]
      \else
        {\Huge\bfseries\color{coulPart} \coursTitre}\\[3pt]
        {\Large\bfseries\color{coulPart} \coursSurTitre}\\[5pt]
      \fi
      {\normalsize \coursSousTitre}%
      \\[2pt]{\small\itshape Version élève}
    \end{tcolorbox}
  \end{center}%
}
\makeatother

\chapitrecours{7}{Pourcentages}
\niveaucours{Seconde}
\titrecours{POURCENTAGES}{Cours --- Seconde}

\begin{document}

\begin{remarquebox}[Comment utiliser ce cours]
Chaque partie commence par ce qu'il faut \textbf{savoir} (définitions, théorèmes), puis
vient ce qu'il faut \textbf{savoir faire} : une méthode, avec des
\textbf{exemples} faits en classe, puis des exercices
\og\textbf{À vous de jouer}\fg{} à chercher seul.

\smallskip
Sur cette feuille, seuls les \textbf{énoncés} des exemples figurent : la résolution,
signalée par la mention \emph{(au tableau)}, est à rédiger dans le cahier.

\end{remarquebox}

\partiecours{Appliquer et calculer un pourcentage}

\begin{definitionbox}[Définition 1 --- Prendre un pourcentage]
Prendre $t\,\%$ d'une quantité, c'est multiplier cette quantité par $\dfrac{t}{100}$.

\smallskip
Par exemple, $20\,\%$ de $45$ : $\dfrac{20}{100}\times 45=0{,}2\times 45=9$.
\end{definitionbox}

\begin{remarquebox}
$100\,\%$ d'une quantité, c'est toute la quantité ; $50\,\%$, c'est la moitié ;
$25\,\%$, le quart ; $200\,\%$, le double.
\end{remarquebox}

\begin{methodebox}[Méthode 1 --- Appliquer un pourcentage]
\etapeexemples
\begin{questions}[label=\alph*)]
  \item Un lycée compte $1\,250$ élèves, dont $36\,\%$ sont demi-pensionnaires.
        Combien d'élèves sont demi-pensionnaires ?
  \item Une console de jeux coûte $380$~€ ; elle a une réduction de
        $15\,\%$. Calculer le montant de la réduction, puis le prix payé.
  \item Calculer $4{,}5\,\%$ de $2\,000$, puis $150\,\%$ de $40$.
\end{questions}

\etape{À vous de jouer}
\begin{questions}[label=\alph*)]
  \item Un village compte $1\,400$ habitants, dont $45\,\%$ ont moins de $30$ ans.
        Combien d'habitants ont moins de $30$ ans ?
  \item Un jean coûte $60$~€ ; il a une réduction de $30\,\%$. Calculer le
        montant de la réduction, puis le prix payé.
  \item Calculer $2{,}5\,\%$ de $800$, puis $120\,\%$ de $35$.
\end{questions}

\end{methodebox}

\afaire{exercices 1 et 2}

\begin{proprietebox}[Propriété 1 --- Quel pourcentage ?]
Dans un groupe de $N$ individus, $n$ ont un caractère commun. Ils représentent
$\dfrac{n}{N}$ du groupe ; écrit en centièmes, ce quotient donne le pourcentage.

\smallskip
Par exemple, $3$ élèves sur $4$ : $\dfrac{3}{4}=0{,}75=\dfrac{75}{100}=75\,\%$.
\end{proprietebox}

\begin{methodebox}[Méthode 2 --- Calculer un pourcentage]
\etapeexemples
\begin{questions}[label=\alph*)]
  \item Dans une classe de $25$ élèves, $9$ font du théâtre. Quel pourcentage des
        élèves font du théâtre ?
  \item Sur un salaire de $1\,600$~€, Inès paie un loyer de $520$~€. Quel
        pourcentage de son salaire représente son loyer ?
  \item Le club de foot A compte $18$ filles sur $40$ licenciés ; le club B,
        $27$ filles sur $75$ licenciés. Quel club a la plus forte part de filles ?
\end{questions}

\etape{À vous de jouer}
\begin{questions}[label=\alph*)]
  \item Dans une classe de $35$ élèves, $21$ viennent au lycée en bus. Quel
        pourcentage des élèves vient en bus ?
  \item Sur un budget mensuel de $1\,850$~€, une famille dépense $296$~€ en
        transports. Quel pourcentage du budget représentent les transports ?
  \item Le lycée A compte $84$ internes sur $600$ élèves ; le lycée B, $117$
        internes sur $900$ élèves. Quel lycée a la plus forte part d'internes ?
\end{questions}

\end{methodebox}

\afaire{exercices 3 à 6}

\partiecours{Proportions}

\begin{definitionbox}[Définition 2 --- Proportion]
\begin{minipage}[c]{0.58\linewidth}\raggedright
La \textbf{population de référence} est le groupe entier, d'effectif $N$ ; une
\textbf{sous-population} en est une partie, d'effectif $n$.

\smallskip
La \textbf{proportion} de la sous-population dans la population de référence est
\[p=\dfrac{n}{N}.\]
C'est un nombre compris entre $0$ et $1$, qu'on écrit en fraction, en écriture
décimale ou en pourcentage.
\end{minipage}\hfill
\begin{minipage}[c]{0.38\linewidth}\centering
\begin{tikzpicture}
  \draw[thick,coulDef] (0,0) rectangle (3.6,2.3);
  \node[font=\footnotesize,anchor=north west] at (0.05,2.3) {population : $N$};
  \draw[thick,fill=coulDef!15] (2,0.95) ellipse (1.2 and 0.6);
  \node[font=\footnotesize,align=center] at (2,0.95) {sous-population\\ $n$};
\end{tikzpicture}
\end{minipage}
\end{definitionbox}

\begin{proprietebox}[Propriété 2 --- Effectif et effectif total]
Quand on connaît la proportion $p$ :
\begin{itemize}
  \item l'effectif de la sous-population est $n=p\times N$ (c'est la méthode 1) ;
  \item l'effectif total $N$ vérifie $p\times N=n$, donc $N=\dfrac{n}{p}$.
\end{itemize}
\end{proprietebox}

\begin{methodebox}[Méthode 3 --- Associer effectif, proportion et effectif total]
\etapeexemples
Un club d'escalade compte $240$ licenciés, dont $45\,\%$ sont mineurs. Parmi les
licenciés, $96$ sont des femmes, et $48$ d'entre elles sont mineures.
\begin{questions}[label=\alph*)]
  \item Combien de licenciés sont mineurs ?
  \item Quelle est la proportion de femmes parmi les licenciés ?
  \item Quelle est la proportion de mineures parmi les femmes ? La comparer à la
        proportion de mineurs dans tout le club.
  \item Un autre club compte $63$ licenciés mineurs, soit $35\,\%$ de ses
        licenciés. Combien ce club a-t-il de licenciés ?
\end{questions}

\etape{À vous de jouer}
Un lycée compte $800$ élèves, dont $30\,\%$ sont externes. Parmi les élèves,
$440$ sont des filles, et $154$ d'entre elles sont externes.
\begin{questions}[label=\alph*)]
  \item Combien d'élèves sont externes ?
  \item Quelle est la proportion de filles parmi les élèves ?
  \item Quelle est la proportion d'externes parmi les filles ? La comparer à la
        proportion d'externes dans tout le lycée.
  \item Un autre lycée compte $91$ latinistes, soit $14\,\%$ de ses élèves.
        Combien ce lycée a-t-il d'élèves ?
\end{questions}

\end{methodebox}

\afaire{exercices 7 à 11}

\partiecours{Proportions échelonnées}

\begin{exemplebox}[Exemple --- Un pourcentage d'un pourcentage]
Dans un car de $50$ passagers, $40\,\%$ sont des scolaires, et $60\,\%$ des
scolaires sont des filles.
\begin{calculs}
\text{scolaires} &= 0{,}4\times 50 &\qquad& \rem{$40\,\%$ des passagers}\\
\text{scolaires} &= 20 &&\\
\text{filles scolaires} &= 0{,}6\times 20 && \rem{$60\,\%$ des scolaires}\\
\text{filles scolaires} &= 12 &&
\end{calculs}
Dans le car, la proportion de filles scolaires est $\dfrac{12}{50}=0{,}24$ ; or
$0{,}6\times 0{,}4=0{,}24$ : il suffisait de multiplier les deux proportions.
\end{exemplebox}

\begin{proprietebox}[Propriété 3 --- Proportions échelonnées]
\begin{minipage}[c]{0.58\linewidth}\raggedright
$A$ est une partie de $B$, et $B$ une partie de $C$.
\begin{itemize}
  \item $p_1$ est la proportion de $A$ dans $B$ ;
  \item $p_2$ est la proportion de $B$ dans $C$.
\end{itemize}
Alors la proportion de $A$ dans $C$ est $p=p_1\times p_2$.
\end{minipage}\hfill
\begin{minipage}[c]{0.38\linewidth}\centering
\begin{tikzpicture}
  \draw[thick,fill=coulProp!8] (0,0) ellipse (1.9 and 1.15);
  \draw[thick,fill=coulProp!20] (0.45,-0.1) ellipse (1.25 and 0.8);
  \draw[thick,fill=coulProp!40] (0.8,-0.15) ellipse (0.6 and 0.45);
  \node at (-1.35,0.2) {$C$};
  \node at (-0.4,0.15) {$B$};
  \node at (0.8,-0.15) {$A$};
\end{tikzpicture}
\end{minipage}
\end{proprietebox}

\demo{On note $n_A$, $n_B$ et $n_C$ les effectifs de $A$, de $B$ et de $C$.}
\begin{calculs}
p &= \dfrac{n_A}{n_C} &\qquad& \rem{proportion de $A$ dans $C$}\\
p &= \dfrac{n_A}{n_B}\times\dfrac{n_B}{n_C} && \rem{en simplifiant par $n_B$, on retrouve $\frac{n_A}{n_C}$}\\
p &= p_1\times p_2 && \rem{$\frac{n_A}{n_B}=p_1$ et $\frac{n_B}{n_C}=p_2$}
\end{calculs}
Conclusion : la proportion de $A$ dans $C$ est le produit $p_1\times p_2$.

\begin{methodebox}[Méthode 4 --- Calculer une proportion échelonnée]
\etapeexemples
\begin{questions}[label=\alph*)]
  \item Dans un lycée, $40\,\%$ des élèves sont inscrits à l'association sportive,
        et $35\,\%$ des inscrits font du volley-ball. Quelle proportion des élèves
        du lycée fait du volley-ball à l'association sportive ?
  \item Ce lycée compte $1\,500$ élèves. Combien font du volley-ball à
        l'association sportive ?
  \item Dans une ville, $50\,\%$ des habitants ont un vélo ; $60\,\%$ d'entre eux
        vont au travail à vélo, et $30\,\%$ de ces derniers ont un vélo électrique.
        Quelle proportion des habitants va au travail en vélo électrique ?
\end{questions}

\etape{À vous de jouer}
\begin{questions}[label=\alph*)]
  \item Dans une entreprise, $60\,\%$ des salariés viennent en train, et $45\,\%$ de
        ceux-ci ont un abonnement annuel. Quelle proportion des salariés vient en
        train avec un abonnement annuel ?
  \item Cette entreprise compte $2\,400$ salariés. Combien viennent en train avec
        un abonnement annuel ?
  \item Dans un lycée, $80\,\%$ des élèves ont un téléphone portable ; $75\,\%$
        d'entre eux ont une application de révision, et $40\,\%$ de ces derniers
        l'utilisent chaque jour. Quelle proportion des élèves utilise chaque jour
        une application de révision ?
\end{questions}

\end{methodebox}

\afaire{exercices 12 à 15}

\partiecours{Coefficient multiplicateur}

Dans la méthode 1, la console à $380$~€, qui a une réduction de $15\,\%$, coûte
$380-57=323$~€. On paie $100\,\%-15\,\%=85\,\%$ du prix : une seule multiplication
suffit, $0{,}85\times 380=323$.

\begin{definitionbox}[Définition 3 --- Taux d'évolution]
Une hausse de $15\,\%$ a pour \textbf{taux d'évolution} $t=+15\,\%=0{,}15$ ; une baisse
de $15\,\%$ a pour taux d'évolution $t=-15\,\%=-0{,}15$.

\smallskip
Le taux d'évolution est positif pour une hausse, négatif pour une baisse.
\end{definitionbox}

\begin{proprietebox}[Propriété 4 --- Coefficient multiplicateur]
\begin{minipage}[c]{0.62\linewidth}\raggedright
Une évolution de taux $t$ multiplie la valeur par $C_m=1+t$, le
\textbf{coefficient multiplicateur} de l'évolution : la valeur de départ $V_0$
devient $V_1=V_0\times C_m$.

\smallskip
Hausse de $15\,\%$ : $C_m=1+0{,}15=1{,}15$ ; baisse de $15\,\%$ :
$C_m=1-0{,}15=0{,}85$.
\end{minipage}\hfill
\begin{minipage}[c]{0.34\linewidth}\centering
\begin{tikzpicture}
  \node (a) at (0,0) {$V_0$};
  \node (b) at (2.6,0) {$V_1$};
  \draw[->,thick,coulProp] (a) to[bend left=40]
    node[above,font=\footnotesize]{$\times\,C_m$} (b);
\end{tikzpicture}
\end{minipage}
\end{proprietebox}

\demo{Pour une baisse, $t$ est négatif : ajouter $t\times V_0$, c'est retirer.}
\begin{calculs}
V_1 &= V_0+t\times V_0 &\qquad& \rem{on ajoute $t\times V_0$ : $15\,\%$ de $V_0$ si $t=0{,}15$ (partie I)}\\
V_1 &= V_0\times(1+t) && \rem{on factorise par $V_0$}
\end{calculs}
Conclusion : une évolution de taux $t$ multiplie la valeur par $1+t$.

\begin{center}
{\renewcommand{\arraystretch}{1.8}\setlength{\tabcolsep}{12pt}
\begin{tabular}{|c|c|c|}
\hline
\textbf{Évolution} & \textbf{Taux $t$} & \textbf{Coefficient $C_m=1+t$}\\
\hline
hausse de $12\,\%$ & $0{,}12$ & $1+0{,}12=1{,}12$\\
\hline
hausse de $5\,\%$ & $0{,}05$ & $1+0{,}05=1{,}05$\\
\hline
baisse de $30\,\%$ & $-0{,}3$ & $1-0{,}3=0{,}7$\\
\hline
baisse de $2\,\%$ & $-0{,}02$ & $1-0{,}02=0{,}98$\\
\hline
\end{tabular}}
\end{center}
Un coefficient $C_m$ plus grand que $1$ traduit une hausse ; plus petit que $1$, une baisse.

\begin{methodebox}[Méthode 5 --- Passer du taux d'évolution au coefficient multiplicateur, et inversement]
\etapeexemples
\begin{questions}[label=\alph*)]
  \item Donner le coefficient multiplicateur $C_m$ d'une hausse de $8\,\%$.
  \item Donner le coefficient multiplicateur $C_m$ d'une baisse de $35\,\%$.
  \item Un prix est multiplié par $1{,}4$. Quel est son taux d'évolution ?
  \item Un prix est multiplié par $0{,}93$. Quel est son taux d'évolution ?
\end{questions}

\etape{À vous de jouer}
\begin{questions}[label=\alph*)]
  \item Donner le coefficient multiplicateur $C_m$ d'une hausse de $7\,\%$.
  \item Donner le coefficient multiplicateur $C_m$ d'une baisse de $22\,\%$.
  \item Un prix est multiplié par $1{,}15$. Quel est son taux d'évolution ?
  \item Un prix est multiplié par $0{,}6$. Quel est son taux d'évolution ?
\end{questions}

\end{methodebox}

\afaire{exercices 16 et 17}

\begin{methodebox}[Méthode 6 --- Calculer une valeur d'arrivée ou de départ]
\etapeexemples
\begin{questions}[label=\alph*)]
  \item La console à $380$~€ a une réduction de $15\,\%$. Calculer le prix payé avec le
        coefficient multiplicateur.
  \item Un loyer de $650$~€ augmente de $3\,\%$. Calculer le nouveau loyer.
  \item Après une hausse de $20\,\%$, un vélo coûte $540$~€. Quel était son prix
        avant la hausse ?
\end{questions}

\etape{À vous de jouer}
\begin{questions}[label=\alph*)]
  \item Un manteau à $240$~€ est soldé de $25\,\%$. Calculer le prix soldé avec le
        coefficient multiplicateur.
  \item Un salaire de $1\,500$~€ augmente de $4\,\%$. Calculer le nouveau salaire.
  \item Après une baisse de $10\,\%$, un casque coûte $63$~€. Quel était son prix
        avant la baisse ?
\end{questions}

\end{methodebox}

\afaire{exercices 18 à 20}

\partiecours{Calculer un taux d'évolution}

\begin{definitionbox}[Définition 4 --- Variation absolue]
Une grandeur passe de la valeur de départ $V_0$ à la valeur d'arrivée $V_1$. Sa
\textbf{variation absolue} est $V_1-V_0$ ; elle a l'unité de la grandeur.
\end{definitionbox}

\begin{proprietebox}[Propriété 5 --- Calculer un taux d'évolution]
Le taux d'évolution (partie IV) de $V_0$ à $V_1$ est
\[t=\dfrac{V_1-V_0}{V_0}.\]
Il n'a pas d'unité et s'écrit souvent en pourcentage ; positif pour une hausse,
négatif pour une baisse.
\end{proprietebox}

\begin{remarquebox}
On a toujours $C_m=1+t$ (propriété 4) : une hausse de $12{,}5\,\%$ a pour coefficient
$1{,}125$ ; une baisse de $15\,\%$, $1+(-0{,}15)=0{,}85$.
\end{remarquebox}

\begin{methodebox}[Méthode 7 --- Calculer un taux d'évolution]
\etapeexemples
Pour chaque grandeur, calculer le taux d'évolution, en pourcentage.
\begin{questions}[label=\alph*)]
  \item Un abonnement de musique en ligne passe de $12$~€ à $13{,}50$~€ par mois.
  \item Un musée reçoit $48\,000$ visiteurs une année, puis $40\,800$ l'année
        suivante.
  \item Une commune passe de $2\,400$ à $2\,530$ habitants (arrondir au dixième de
        pourcentage).
\end{questions}

\etape{À vous de jouer}
Pour chaque grandeur, calculer le taux d'évolution, en pourcentage.
\begin{questions}[label=\alph*)]
  \item Une place de cinéma passe de $15$~€ à $16{,}80$~€.
  \item Une ligne de bus transporte $25\,000$ voyageurs un mois, puis $21\,000$ le
        mois suivant.
  \item Un club passe de $1\,800$ à $1\,950$ adhérents (arrondir au dixième de
        pourcentage).
\end{questions}

\end{methodebox}

\afaire{exercices 21 à 26}

\partiecours{Évolutions successives}

Un article à $100$~€ augmente de $20\,\%$ : il coûte $120$~€. Puis il baisse de
$20\,\%$ : il coûte $120\times 0{,}8=96$~€, et non $100$~€. Des pourcentages
successifs ne s'ajoutent pas.

\begin{proprietebox}[Propriété 6 --- Évolutions successives]
\begin{minipage}[c]{0.52\linewidth}\raggedright
Quand une valeur subit deux évolutions successives, de coefficients $C_{m1}$
puis $C_{m2}$, le \textbf{coefficient multiplicateur global} est le produit
\[C_m=C_{m1}\times C_{m2}.\]
De même pour trois évolutions ou plus.
\end{minipage}\hfill
\begin{minipage}[c]{0.44\linewidth}\centering
\begin{tikzpicture}
  \node (a) at (0,0) {$V_0$};
  \node (b) at (2.2,0) {$V_1$};
  \node (c) at (4.4,0) {$V_2$};
  \draw[->,thick,coulProp] (a) to[bend left=40]
    node[above,font=\footnotesize]{$\times\,C_{m1}$} (b);
  \draw[->,thick,coulProp] (b) to[bend left=40]
    node[above,font=\footnotesize]{$\times\,C_{m2}$} (c);
  \draw[->,thick,coulExo] (a) to[bend right=30]
    node[below,font=\footnotesize]{$\times\,C_{m1}\times C_{m2}$} (c);
\end{tikzpicture}
\end{minipage}
\end{proprietebox}

\demo{}
\begin{calculs}
V_2 &= V_1\times C_{m2} &\qquad& \rem{seconde évolution}\\
V_2 &= V_0\times C_{m1}\times C_{m2} && \rem{première évolution : $V_1=V_0\times C_{m1}$}
\end{calculs}
Conclusion : passer de $V_0$ à $V_2$, c'est multiplier par
$C_{m1}\times C_{m2}$.

\begin{methodebox}[Méthode 8 --- Calculer un taux d'évolution global]
\etapeexemples
\begin{questions}[label=\alph*)]
  \item Un prix augmente de $25\,\%$, puis baisse de $12\,\%$. Quelle est son
        évolution globale ?
  \item Le nombre d'abonnés d'un journal baisse de $10\,\%$ par an pendant trois
        ans. Quelle est son évolution globale ?
  \item Un plat à $80$~€ augmente de $5\,\%$ deux années de suite. Quelle est son
        évolution globale ? Quel est son prix final ?
\end{questions}

\etape{À vous de jouer}
\begin{questions}[label=\alph*)]
  \item Un prix augmente de $50\,\%$, puis baisse de $20\,\%$. Quelle est son
        évolution globale ?
  \item La valeur d'une voiture baisse de $20\,\%$ par an pendant trois ans.
        Quelle est son évolution globale ?
  \item Une carte de transport à $60$~€ augmente de $10\,\%$ deux années de suite.
        Quelle est son évolution globale ? Quel est son prix final ?
\end{questions}

\end{methodebox}

\afaire{exercices 27 à 31}

\partiecours{Évolution réciproque}

Dans la partie VI, une baisse de $20\,\%$ n'a pas compensé une hausse de $20\,\%$.
Quelle évolution ramène une valeur à son point de départ ?

\begin{definitionbox}[Définition 5 --- Évolution réciproque]
Une valeur passe de $V_0$ à $V_1$. L'\textbf{évolution réciproque} est celle qui
ramène $V_1$ à $V_0$.
\end{definitionbox}

\begin{proprietebox}[Propriété 7 --- Coefficient de l'évolution réciproque]
\begin{minipage}[c]{0.58\linewidth}\raggedright
Le coefficient multiplicateur de l'évolution réciproque est l'inverse de celui de
l'évolution de départ :
\[C'_m=\dfrac{1}{C_m}.\]
\end{minipage}\hfill
\begin{minipage}[c]{0.38\linewidth}\centering
\begin{tikzpicture}
  \node (a) at (0,0) {$V_0$};
  \node (b) at (2.6,0) {$V_1$};
  \draw[->,thick,coulProp] (a) to[bend left=40]
    node[above,font=\footnotesize]{$\times\,C_m$} (b);
  \draw[->,thick,coulExo] (b) to[bend left=40]
    node[below,font=\footnotesize]{$\times\,\dfrac{1}{C_m}$} (a);
\end{tikzpicture}
\end{minipage}
\end{proprietebox}

\demo{}
\begin{calculs}
V_0\times C_m &= V_1 &\qquad& \rem{l'évolution de départ}\\
V_0 &= \dfrac{V_1}{C_m} && \rem{on divise par $C_m$}\\
V_0 &= V_1\times\dfrac{1}{C_m} &&
\end{calculs}
Conclusion : pour revenir de $V_1$ à $V_0$, on multiplie par $\dfrac{1}{C_m}$.

\begin{methodebox}[Méthode 9 --- Calculer un taux d'évolution réciproque]
\etapeexemples
Pour chaque évolution, quelle évolution ramène la valeur à son point de départ ?
\begin{questions}[label=\alph*)]
  \item Les ventes d'un magasin baissent de $20\,\%$.
  \item Le prix d'un appartement augmente de $60\,\%$.
  \item Le prix du kilo de cerises augmente de $15\,\%$ (arrondir au dixième de
        pourcentage).
\end{questions}

\etape{À vous de jouer}
Pour chaque évolution, quelle évolution ramène la valeur à son point de départ ?
\begin{questions}[label=\alph*)]
  \item Le nombre de clients d'un restaurant baisse de $50\,\%$.
  \item Le prix d'un forfait de ski augmente de $25\,\%$.
  \item Le prix d'un ordinateur baisse de $30\,\%$ (arrondir au dixième de
        pourcentage).
\end{questions}

\end{methodebox}

\afaire{exercices 32 à 36}

\end{document}
